%% TOAN12_C1_ONTAP1.tex -- Đề ôn tập chương I, đề số 1
%% Câu hỏi lấy từ ngân hàng đề (id ghi trong chú thích). Đáp án: \dapan{khoá}{hướng dẫn}.
\begin{deontap}{1}
\phanI
%% C1-B1-D05-NB-008
\begin{cau}Cho hàm số $y=f(x)$ có bảng xét dấu của đạo hàm như sau:
\begin{center}
\begin{tikzpicture}
\tkzTabInit[lgt=1.4,espcl=2.0]{$x$/0.7,$f'(x)$/0.7}{$-\infty$,$-3$,$4$,$+\infty$}
\tkzTabLine{ ,-,z,+,z,-, }
\end{tikzpicture}
\end{center}
Hàm số đã cho đạt cực đại tại điểm nào dưới đây?
\pa{$x=4$.}{$x=-3$.}{$x=5$.}{$x=-4$.}
\dapan{A}{$f'(x)$ đổi dấu từ dương sang âm khi qua $x=4$ nên hàm số đạt cực đại tại đó.}\end{cau}
%% C1-B2-D03-NB-019
\begin{cau}Xét hàm số $y=x+\dfrac{1}{x}$ trên khoảng $(-\infty;0)$. Mệnh đề nào dưới đây \textbf{đúng}?
\pa{Hàm số không có giá trị nhỏ nhất trên khoảng $(-\infty;0)$.}{Hàm số có giá trị nhỏ nhất bằng $-2$.}{Hàm số có giá trị nhỏ nhất bằng $2$.}{Hàm số không có giá trị lớn nhất trên khoảng $(-\infty;0)$.}
\dapan{A}{Ta có $y'=\dfrac{(x - 1) (x + 1)}{x^{2}}$. $y'=0\Leftrightarrow x=-1 \vee x=1$ (loại $x=1$ vì không thuộc $(-\infty;0)$). Xét dấu: $y'>0$ trên $(-\infty;-1)$, $y'<0$ trên $(-1;0)$. Ta có $y(-1)=-2$, $\lim\limits_{x\to-\infty}y=-\infty$, $\lim\limits_{x\to{0}^{-}}y=-\infty$. Do đó $\max y=-2$ (tại $x=-1$); hàm số không có giá trị nhỏ nhất.}\end{cau}
%% C1-B3-D02-NB-007
\begin{cau}Tiệm cận đứng của đồ thị hàm số $y=\dfrac{4x+4}{x-3}$ là đường thẳng
\pa{$x=3$.}{$x=-3$.}{$y=3$.}{$x=4$.}
\dapan{A}{Xét hàm số $y=\dfrac{4x+4}{x-3}$. $\lim\limits_{x\to {3}^{+}}y=+\infty$, $\lim\limits_{x\to {3}^{-}}y=-\infty$ nên $x=3$ là tiệm cận đứng; $\lim\limits_{x\to +\infty}y=4$; $\lim\limits_{x\to -\infty}y=4$. Vậy đồ thị có 1 tiệm cận đứng ($x=3$) và 1 tiệm cận ngang ($y=4$).}\end{cau}
%% C1-B4-D04-NB-029
\begin{cau}Cho hàm số $y=f(x)$ xác định trên $\mathbb{R}\setminus\{-2\}$ và có bảng biến thiên như sau:
\begin{center}
\begin{tikzpicture}
\tkzTabInit[lgt=1.2,espcl=1.9]{$x$/0.7,$y'$/0.7,$y$/2}{$-\infty$,$-3$,$-2$,$-1$,$+\infty$}
\tkzTabLine{ ,-,z,+,d,+,z,-, }
\tkzTabVar{+/$+\infty$,-/$6$,+D-/$+\infty$/$-\infty$,+/$2$,-/$-\infty$}
\end{tikzpicture}
\end{center}
Số nghiệm thực của phương trình $f(x)=7$ là
\pa{$1$.}{$2$.}{$3$.}{$4$.}
\dapan{B}{Số nghiệm của phương trình $f(x)=7$ bằng số giao điểm của đồ thị $y=f(x)$ với đường thẳng $y=7$. Dựa vào bảng biến thiên ta được $2$ nghiệm.}\end{cau}
%% C1-B5-D02-NB-001
\begin{cau}Tổng chi phí (nghìn đồng) để sản xuất $x$ sản phẩm là $C(x)=\dfrac{1}{2}x^{2}+30x+5000$ $(x>0)$. Chi phí trung bình cho mỗi sản phẩm là $\overline{C}(x)=\dfrac{C(x)}{x}$. Chi phí biên khi sản xuất $50$ sản phẩm bằng
\pa{$79\,\text{ nghìn đồng}$.}{$80\,\text{ nghìn đồng}$.}{$78\,\text{ nghìn đồng}$.}{$40\,\text{ nghìn đồng}$.}
\dapan{B}{$C'(x)=x+30$ nên $C'(50)=80$.}\end{cau}
%% C1-B1-D06-TH-044
\begin{cau}Cho hàm số $y=\dfrac{x^{2}}{x-2}$. Đường thẳng đi qua hai điểm cực trị của đồ thị hàm số đã cho đi qua điểm nào dưới đây?
\pa{$M(-3;-8)$.}{$M(-3;-7)$.}{$M(-3;-4)$.}{$M(-3;-6)$.}
\dapan{D}{Tập xác định: $D=\mathbb{R}\setminus\{2\}$. Ta có $y'=\dfrac{x (x - 4)}{(x - 2)^{2}}$; $y'=0\Leftrightarrow x=0 \vee x=4$. Xét dấu: $y'>0$ trên $(-\infty;0)$, $y'<0$ trên $(0;2)$, $y'<0$ trên $(2;4)$, $y'>0$ trên $(4;+\infty)$. Do đó hàm số đạt cực đại tại $x=0$, $y_{\text{CĐ}}=0$; cực tiểu tại $x=4$, $y_{\text{CT}}=8$. Hai điểm cực trị của đồ thị: $A(0;0)$, $B(4;8)$. Đường thẳng $AB$: $y=2x$. Với $x=-3$ thì $y=-6$, nên đáp án là $M(-3;-6)$.}\end{cau}
%% C1-B2-D04-TH-032
\begin{cau}Một bể chứa $200$ lít nước tinh khiết. Người ta bơm vào bể nước muối có nồng độ $25$ gam/lít với tốc độ $20$ lít/phút. Nồng độ muối trong bể sau $t$ phút là $C(t)=\dfrac{500t}{200+20t}$ (gam/lít). Lượng muối (gam) trong bể sau $40$ phút bằng
\pa{$15000\,\text{ g}$.}{$20000\,\text{ g}$.}{$30000\,\text{ g}$.}{$40000\,\text{ g}$.}
\dapan{B}{Lượng muối $=25\cdot20\cdot40=20000$ (gam).}\end{cau}
%% C1-B3-D01-TH-035
\begin{cau}Cho hàm số $y=f(x)$ xác định trên $\mathbb{R}\setminus\{3\}$, liên tục trên mỗi khoảng xác định và có bảng biến thiên như sau:
\begin{center}
\begin{tikzpicture}
\tkzTabInit[lgt=1.2,espcl=1.9]{$x$/0.7,$y'$/0.7,$y$/2}{$-\infty$,$2$,$3$,$+\infty$}
\tkzTabLine{ ,+,z,-,d,+, }
\tkzTabVar{-/$-\infty$,+/$0$,-D-/$-\infty$/$-\infty$,+/$+\infty$}
\end{tikzpicture}
\end{center}
Tổng số tiệm cận đứng và tiệm cận ngang của đồ thị hàm số đã cho là
\pa{$1$.}{$2$.}{$3$.}{$0$.}
\dapan{A}{Từ bảng biến thiên: $\lim\limits_{x\to-\infty}f(x)=-\infty$; $\lim\limits_{x\to {3}^{-}}f(x)=-\infty$, $\lim\limits_{x\to {3}^{+}}f(x)=-\infty$; $\lim\limits_{x\to+\infty}f(x)=+\infty$. Do đó đồ thị có 1 tiệm cận đứng ($x=3$) và 0 tiệm cận ngang.}\end{cau}
%% C1-B4-D02-TH-017
\begin{cau}Bảng biến thiên dưới đây là của hàm số nào?
\begin{center}
\begin{tikzpicture}
\tkzTabInit[lgt=1.2,espcl=2.6]{$x$/0.7,$y'$/0.7,$y$/2}{$-\infty$,$-2$,$+\infty$}
\tkzTabLine{ ,+,d,+, }
\tkzTabVar{-/$2$,+D-/$+\infty$/$-\infty$,+/$2$}
\end{tikzpicture}
\end{center}
\pa{$y=\dfrac{2x+2}{x+2}$.}{$y=\dfrac{2x+2}{x-2}$.}{$y=\dfrac{-2x+2}{x+2}$.}{$y=\dfrac{2x+5}{x+2}$.}
\dapan{A}{Từ bảng biến thiên: tiệm cận đứng $x=-2$, tiệm cận ngang $y=2$, hàm số đồng biến trên từng khoảng xác định. Chỉ hàm số $y=\dfrac{2x+2}{x+2}$ thỏa mãn (các phương án còn lại sai tiệm cận hoặc sai chiều biến thiên).}\end{cau}
%% C1-B1-D03-VD-017
\begin{cau}Điện lượng truyền qua tiết diện của một dây dẫn tính đến thời điểm $t$ (giây) là $Q(t)=t^{3}-9t^{2}+5t$ (culông). Cường độ dòng điện tức thời là $I(t)=Q'(t)$ (ampe). Cường độ dòng điện tại thời điểm $t=2$ giây bằng
\pa{$-19\,\text{ A}$.}{$-18\,\text{ A}$.}{$-17\,\text{ A}$.}{$-20\,\text{ A}$.}
\dapan{A}{Ta có $I(t)=Q'(t)=3t^{2}-18t+5$. $I(2)=-19$.}\end{cau}
%% C1-B2-D02-VD-023
\begin{cau}Một quả bóng được ném thẳng đứng lên trên từ độ cao $1$ m với vận tốc ban đầu $12$ m/s. Độ cao của quả bóng (tính bằng mét) sau $t$ giây là $h(t)=1+12t-5t^2$ (bỏ qua sức cản không khí, lấy $g=10$ m/s$^2$). Thời điểm quả bóng đạt độ cao lớn nhất là
\hinh{\begin{tikzpicture}[scale=0.9,font=\footnotesize,line join=round,line cap=round,>=Stealth]
\draw[thick] (-1.6,0)--(1.8,0); \fill[pattern=north east lines] (-1.6,-.15) rectangle (1.8,0);
\fill[gray!30] (-.25,0) rectangle (.25,.9); \draw (-.25,0) rectangle (.25,.9);
\draw[dashed] (0,.9)--(0,2.6); \shade[ball color=gray!40] (0,2.75) circle (.14);
\draw[->,thick] (.35,1.1)--(.35,2) node[midway,right]{$v_0$};
\draw[<->] (-.55,0)--(-.55,.9) node[midway,left]{$h_0$};
\draw[->] (-1.2,0)--(-1.2,3) node[above]{$h$};
\end{tikzpicture}}

\pa{$\dfrac{1}{5}\,\mathrm{s}$.}{$\dfrac{6}{5}\,\mathrm{s}$.}{$\dfrac{5}{6}\,\mathrm{s}$.}{$\dfrac{3}{5}\,\mathrm{s}$.}
\dapan{B}{Ta có $v(t)=h'(t)=12-10t$. $h'(t)=0\Leftrightarrow t=\dfrac{6}{5}$.}\end{cau}
%% C1-B5-D04-VD-026
\begin{cau}Một xưởng in áo có chi phí cố định là $1000$ nghìn đồng và chi phí in mỗi chiếc áo là $50$ nghìn đồng. Chi phí trung bình để in $x$ chiếc áo là $\overline{C}(x)=\dfrac{50x+1000}{x}$ (nghìn đồng/chiếc), $x>0$. Chi phí cố định bằng
\pa{$2000\,\text{ nghìn đồng}$.}{$1000\,\text{ nghìn đồng}$.}{$1500\,\text{ nghìn đồng}$.}{$500\,\text{ nghìn đồng}$.}
\dapan{B}{Chi phí cố định là $1000$ nghìn đồng (tử số khi $x=0$).}\end{cau}
\phanII
%% C1-B1-D04-NB-038
\begin{cau}Cho hàm số $y=-2x^{4}+4x^{2}+2$. Xét tính đúng sai của các mệnh đề sau:
\begin{dungsai}
  \y Điểm cực đại của đồ thị hàm số là $(4;1)$. \oDS
  \y Hàm số đạt cực đại tại $x=1$. \oDS
  \y Đạo hàm của hàm số là $y'=- 8 x (x - 1) (x + 1)$. \oDS
  \y Điểm cực đại của đồ thị hàm số là $(4;-1)$. \oDS
\end{dungsai}
\dapan{S,Đ,Đ,S}{Tập xác định: $D=\mathbb{R}$. Ta có $y'=- 8 x (x - 1) (x + 1)$; $y'=0\Leftrightarrow x=-1 \vee x=0 \vee x=1$. Xét dấu: $y'>0$ trên $(-\infty;-1)$, $y'<0$ trên $(-1;0)$, $y'>0$ trên $(0;1)$, $y'<0$ trên $(1;+\infty)$. Do đó hàm số đạt cực đại tại $x=-1$, $y_{\text{CĐ}}=4$; cực tiểu tại $x=0$, $y_{\text{CT}}=2$; cực đại tại $x=1$, $y_{\text{CĐ}}=4$. a) Sai: điểm cực đại là $(1;4)$. b) Đúng. c) Đúng. d) Sai: điểm cực đại là $(-1;4)$.}\end{cau}
%% C1-B2-D04-TH-061
\begin{cau}Từ một tấm tôn hình chữ nhật có chiều rộng $36$ cm, người ta gập hai mép (mỗi mép rộng $x$ cm) vuông góc lên để làm một cái máng xối có mặt cắt ngang là hình chữ nhật. Xét tính đúng sai của các mệnh đề sau:
\begin{dungsai}
  \y Diện tích mặt cắt ngang của máng được tính bởi hàm số $f(x)=x(36-2x)$. \oDS
  \y Diện tích mặt cắt ngang của máng đạt giá trị lớn nhất khi $x=18$. \oDS
  \y Giá trị lớn nhất của diện tích mặt cắt ngang của máng bằng $162$. \oDS
  \y Điều kiện của $x$ là $x\in(0;18)$. \oDS
\end{dungsai}
\dapan{Đ,S,Đ,Đ}{Mặt cắt ngang là hình chữ nhật có đáy $36-2x$ và chiều cao $x$, diện tích $S(x)=x(36-2x)$ với $0<x<18$. Xét $f(x)=x(36-2x)$, ta có $f'(x)=- 4 (x - 9)$, $f'(x)=0$ có nghiệm thuộc miền xác định là $x=9$; $f'$ đổi dấu từ dương sang âm qua nghiệm này. Do đó diện tích mặt cắt ngang của máng lớn nhất bằng $162$ khi $x=9$. a) Đúng. b) Sai: giá trị tối ưu đạt tại $x=9$. c) Đúng. d) Đúng.}\end{cau}
%% C1-B3-D03-TH-064
\begin{cau}Cho hàm số $y=-x+1+\dfrac{2}{x-1}$. Xét tính đúng sai của các mệnh đề sau:
\begin{dungsai}
  \y Đồ thị hàm số có tất cả $3$ đường tiệm cận. \oDS
  \y Đường thẳng $y=-x+3$ là tiệm cận xiên của đồ thị hàm số. \oDS
  \y Đồ thị hàm số có tất cả $2$ đường tiệm cận. \oDS
  \y Đường thẳng $y=-x$ là tiệm cận xiên của đồ thị hàm số. \oDS
\end{dungsai}
\dapan{S,S,Đ,S}{Xét hàm số $y=-x+1+\dfrac{2}{x-1}$. Chia đa thức: $y=-x+1+\dfrac{2}{x - 1}$; $a=\lim\limits_{x\to +\infty}\dfrac{y}{x}=-1$, $b=\lim\limits_{x\to +\infty}(y+x)=1$ nên $y=-x+1$ là tiệm cận xiên (khi $x\to +\infty$); $a=\lim\limits_{x\to -\infty}\dfrac{y}{x}=-1$, $b=\lim\limits_{x\to -\infty}(y+x)=1$ nên $y=-x+1$ là tiệm cận xiên (khi $x\to -\infty$); $\lim\limits_{x\to {1}^{+}}y=+\infty$ nên $x=1$ là tiệm cận đứng. a) Sai: có $2$ đường tiệm cận. b) Sai. c) Đúng. d) Sai: tiệm cận xiên là $y=-x+1$.}\end{cau}
%% C1-B4-D04-VD-038
\begin{cau}Một vật dao động điều hòa có phương trình li độ $x(t)=2\cos(\pi t)$ (cm), với $t$ tính bằng giây. Xét tính đúng sai của các mệnh đề sau:
\hinh{\begin{tikzpicture}[scale=0.9,font=\footnotesize,line join=round,line cap=round,>=Stealth]
\draw[thick] (0,0)--(0,1.2); \fill[pattern=north east lines] (-.2,0) rectangle (0,1.2);
\draw[decorate,decoration={coil,segment length=5pt,amplitude=4pt}] (0,.6)--(2,.6); \filldraw[fill=gray!40] (2,.35) rectangle (2.5,.85);
\draw[->] (2.25,0)--(3.5,0) node[right]{$x$}; \draw[dashed] (2.25,0)--(2.25,.35);
\end{tikzpicture}}

\begin{dungsai}
  \y Tốc độ cực đại của vật là $\pi\,\mathrm{cm/s}$. \oDS
  \y Li độ của vật tại thời điểm $t=\dfrac{1}{6}$ giây là $\dfrac{3 \sqrt{3}}{2}\,\mathrm{cm}$. \oDS
  \y Vận tốc của vật tại thời điểm $t=\dfrac{1}{6}$ giây là $- \pi\,\mathrm{cm/s}$. \oDS
  \y Gia tốc của vật tại thời điểm $t=\dfrac{1}{6}$ giây là $- \sqrt{3} \pi^{2}\,\mathrm{cm/s^2}$. \oDS
\end{dungsai}
\dapan{S,S,Đ,Đ}{Ta có $x(\dfrac{1}{6})=\sqrt{3}$. Vận tốc $v(t)=x'(t)=-2 \pi\sin(\pi t)$ (cm/s), gia tốc $a(t)=v'(t)=-2 \pi^{2}\cos(\pi t)$ (cm/s$^2$). $a(\dfrac{1}{6})=- \sqrt{3} \pi^{2}$. Vận tốc $v(t)=x'(t)=-2 \pi\sin(\pi t)$ (cm/s), gia tốc $a(t)=v'(t)=-2 \pi^{2}\cos(\pi t)$ (cm/s$^2$). Vì $|\sin(\pi t)|\le1$ nên $|v(t)|\le 2 \pi$, dấu bằng xảy ra khi $\sin(\pi t)=\pm1$. Vận tốc $v(t)=x'(t)=-2 \pi\sin(\pi t)$ (cm/s), gia tốc $a(t)=v'(t)=-2 \pi^{2}\cos(\pi t)$ (cm/s$^2$). $v(\dfrac{1}{6})=- \pi$. a) Sai: giá trị đúng là $2 \pi\,\mathrm{cm/s}$. b) Sai: giá trị đúng là $\sqrt{3}\,\mathrm{cm}$. c) Đúng. d) Đúng.}\end{cau}
\phanIII
%% C1-B1-D01-TH-080
\begin{cau}Cho hàm số $y=f(x)$ có đồ thị hàm số $y=f'(x)$ như hình vẽ.
\begin{center}
\begin{tikzpicture}[xscale=0.69,yscale=0.8,>=stealth]
\begin{scope}
\clip (-4.70,-3.30) rectangle (5.70,3.30);
\draw[thick,smooth,samples=200,domain=-4.70:5.70] plot(\x,{(((0.195918)*\x+(-0.195918))*\x+(-2.351020))});
\end{scope}
\draw[->] (-4.70,0)--(5.70,0) node[below]{$x$};
\draw[->] (0,-3.40)--(0,3.50) node[left]{$y$};
\node[below left] at (0,0){$O$};
\node[below left] at (-3.00,0){$-3$};
\node[below right] at (4.00,0){$4$};
\end{tikzpicture}
\end{center}
Hàm số $f(x)$ nghịch biến trên đúng một khoảng bị chặn. Tính tổng các số nguyên thuộc khoảng đó.\otraloi
\dapan{3}{Xét dấu $f'(x)$: $f'(x)>0$ trên $(-\infty;-3)$ và $(4;+\infty)$; $f'(x)<0$ trên $(-3;4)$. Khoảng nghịch biến bị chặn duy nhất của hàm số là $(-3;4)$; các số nguyên thuộc khoảng này là -2, -1, 0, 1, 2, 3. Kết quả là $3$.}\end{cau}
%% C1-B2-D04-TH-069
\begin{cau}Từ một tấm tôn hình chữ nhật có chiều rộng $60$ cm, người ta gập hai mép (mỗi mép rộng $x$ cm) vuông góc lên để làm một cái máng xối có mặt cắt ngang là hình chữ nhật. Tìm $x$ (đơn vị: cm) để diện tích mặt cắt ngang của máng lớn nhất.\otraloi
\dapan{15}{Mặt cắt ngang là hình chữ nhật có đáy $60-2x$ và chiều cao $x$, diện tích $S(x)=x(60-2x)$ với $0<x<30$. Xét $f(x)=x(60-2x)$, ta có $f'(x)=- 4 (x - 15)$, $f'(x)=0$ có nghiệm thuộc miền xác định là $x=15$; $f'$ đổi dấu từ dương sang âm qua nghiệm này. Do đó diện tích mặt cắt ngang của máng lớn nhất bằng $450$ khi $x=15$. Kết quả: $15$.}\end{cau}
%% C1-B1-D05-VD-059
\begin{cau}Cho hàm số $y=f(x)$ có đạo hàm $f'(x)=-x(x-4)(x^2+4)(x+3)(x+1)$ với mọi $x\in\mathbb{R}$. Tính tổng tất cả các điểm cực trị của hàm số $g(x)=f(x^{2})$.\otraloi
\dapan{0}{Ta có $f'(x)=0\Leftrightarrow x=-3 \vee x=-1 \vee x=0 \vee x=4$. Bảng xét dấu cho thấy $f'(x)$ đổi dấu tại $x=-3$ (từ âm sang dương), $x=-1$ (từ dương sang âm), $x=0$ (từ âm sang dương), $x=4$ (từ dương sang âm). Ta có $g'(x)=2 x\cdot f'(x^{2})$. Xét dấu: $g'(x)>0$ trên $(-\infty;-2)$ và $(0;2)$; $g'(x)<0$ trên $(-2;0)$ và $(2;+\infty)$. Các điểm cực trị của $g$: $x=-2$, $x=0$, $x=2$. Vậy kết quả là $0$.}\end{cau}
%% C1-B3-D04-VD-056
\begin{cau}Cho hàm số $y=\dfrac{2x^{2}+x+3}{x+1}$. Khoảng cách từ gốc tọa độ đến giao điểm của hai đường tiệm cận bằng (làm tròn kết quả đến hàng phần trăm).\otraloi
\dapan{3,16}{Đồ thị hàm số $y=\dfrac{2x^{2}+x+3}{x+1}$ có tiệm cận đứng $x=-1$ và tiệm cận xiên $y=2x-1$. Giao điểm $I(-1;-3)$, $OI=\sqrt{10}$. Kết quả là $\sqrt{10}\approx 3,16$.}\end{cau}
%% C1-B4-D01-VD-052
\begin{cau}Một chất điểm chuyển động thẳng có phương trình $x(t)=t^{3}-12t^{2}+45t$ (mét), với $t$ (giây), $0\le t\le 7$. Tìm thời điểm đầu tiên chất điểm dừng lại (đổi chiều chuyển động) (đơn vị: $\mathrm{s}$).\otraloi
\dapan{3}{Vận tốc $v(t)=x'(t)=3t^{2}-24t+45=3(t-3)(t-5)$. $v(t)=0\Leftrightarrow t=3$ hoặc $t=5$; chất điểm dừng lần đầu tại $t=3$. Kết quả: $3$.}\end{cau}
%% C1-B5-D03-VD-050
\begin{cau}Nhiệt độ ($^\circ$C) tại một thành phố $t$ giờ sau nửa đêm ($0\le t\le 24$) được mô hình hoá bởi $T(t)=-\dfrac{20}{729}t^{3}+\dfrac{190}{243}t^{2}-\dfrac{1400}{243}t+\dfrac{22372}{729}$. Tìm độ dài khoảng thời gian (giờ) nhiệt độ tăng liên tục.\otraloi
\dapan{9}{$T'(t)=-\dfrac{20}{243}t^{2}+\dfrac{380}{243}t-\dfrac{1400}{243}$; $T'(t)=0\Leftrightarrow t=5$ hoặc $t=14$; nhiệt độ giảm trên $(0;5)$, tăng trên $(5;14)$, giảm trên $(14;24)$. Kết quả: $9$.}\end{cau}
\end{deontap}
